\section{A summable disagreement exploration on the evolving old core}
\label{erc:section}

We couple two retained components whose roots differ by one leaf. The proof charges the future families born from affected cores and sums the resulting disagreements over the evolving core tree.

Write
\[
 W_T(s)=\sup_x\mathbb E_x\int_0^T e^{-st}C_{\rm tag}(t)\,dt,
 \qquad
 V_{s,T}(y)=\int_0^T e^{-st}P_t^\infty C(y)\,dt.
\]
The first supremum is over the admissible newborn-environment
class, including finite inert-boundary profiles.  A future child in a
finite component with an inert ancestor belongs to this class.
For finite $T$, core-count Yule domination gives
$W_T(s)\le\int_0^T e^{(1-s)t}\,dt<\infty$.
No finiteness of $W(s)=\sup_TW_T(s)$ will be assumed in the comparison.

There is no enlargement of the supremum hidden in the use of a family
inside a finite inert tree.  Replace its ancestor's inert stub by an
edge to a genuine parent of degree $R$, completed by $R-1$ leaves.
The full degree at the old stub endpoint is unchanged.  Couple all
interior transitions, treating outward choices as null in the inert
copy.  The interiors agree until a parent-to-interior arrival, whose
rate is at most $1/R$ because parent degrees never decrease.  Thus the
probability of disagreement through a fixed horizon $T$ is at most
$T/R$.

The tagged newborn core counts have Yule second moments bounded
uniformly in $R$.  Cauchy--Schwarz therefore bounds the discrepancy in
their expected finite-horizon discounted rewards by a constant
depending only on $s,T$ times $\sqrt{T/R}$, which tends to zero.
Each genuine finite exterior has reward at most $W_T(s)$, so the same
holds in the inert limit.  This argument applies to the fixed finite
state encountered at any charged birth.  It uses only an inert
boundary: it does not identify a constant-rate controlled environment
with an original finite exterior.
If a finite exterior profile has several inert stubs, complete each
one in this way.  For $b$ stubs the disagreement probability is at most
$bT/R$, and the same fixed-profile limit applies.

\subsection{The coupling uses aggregate transitions}

For a core $u$, let $D_u$ be its full degree, $L_u$ its leaf count,
and $c_u^\partial$ its number of inert stubs.  Here there is one inert
stub at the retained root and no others.  The inert core generator has
the following two transition types at $u$:
\begin{align}
 a_u&=L_u+\sum_{v\sim u,\ v\text{ core}}\frac1{D_v},
 &&\text{add one leaf at }u,\label{erc:gain-rate}\\
 b_u&=\frac{L_u}{D_u}\le1,
 &&\text{promote a leaf of }u\text{ to a new core with one leaf}.
 \label{erc:bud-rate}
\end{align}
The second transition changes $L_u$ and the number of core neighbours
by opposite amounts and leaves $D_u$ unchanged.

Couple two copies by maximal coupling of the aggregate leaf-gain
channels at each paired core, and, separately, by maximal coupling of
their promotion channels.  These are distinct continuous-time channels.
A promotion at $u$ and a leaf gain at a neighbour therefore do not
occur simultaneously in this coupling.  This construction has precisely
the required generator in each marginal.  It does not impose a common
source-clock ring between the copies; the source identity is absent
from the core-state transition.  Leaves attached to a given core are
exchangeable, so paired promotions may always be matched as new paired
core children.

Start the copies from $y$ and $y^+$, where $y^+$ has one additional
leaf at the retained root.  Mark that root infected at time zero.
The word ``infected'' is solely a coupling mark.  It never changes
either marginal process.

Orient the core tree away from the retained root.  Common cores born
before their parent's first infection remain eligible for exploration;
all cores present initially are eligible.
When an eligible uninfected core first receives an unmatched leaf gain,
mark it infected permanently.  At its infection, freeze its list of
currently existing common core neighbours.  Every paired or unpaired
new core child born from it \emph{after} its infection is instead
charged as an entire newborn family, permanently excluded from the
exploration.  The actual evolution and feedback of that family remain
present in the coupling.
Every later descendant of a charged family remains in that same
charged family and is also excluded from the exploration.

An eligible uninfected core has identical local state in both copies.
Its only possible infected neighbour is its predecessor toward the
root.  Indeed, the infected part of the eligible common tree is
connected, and a charged family attaches to that tree only at its
already infected parent.  Before an uninfected core is reached, its
common promotions can grow arbitrary common fringes; those fringes
are included when their attachment core is eventually infected.
Thus no initial or subsequently grown common fringe is discarded.

\subsection{A local reservoir controls all return influences}

Fix an infected eligible core $u$ and reset its age to zero at infection.
Let $d\ge2$ be its smaller full degree then.  The two degrees differ
by exactly one at that time.  Define, at subsequent age $t$,
\[
 Z_u(t)=|D_u^+(t)-D_u^-(t)|,\qquad
 R_u(t)=\#\{\text{union promotions at }u\text{ after infection}\},
\]
and let $B_u(t)$ count its frozen outward neighbours subsequently
infected from $u$.  A simultaneous promotion in the two marginals
counts once in $R_u$, even though two newborn families are charged.
Initially $(Z_u,R_u,B_u)=(1,0,0)$.

There are at most $d-1$ frozen outward candidates.  For the initial
root this follows from its inert stub; for all later infected cores
one neighbour is their infected predecessor.  This is the reason the
estimate below is a root-boundary increment estimate.

Put $h(D)=1-1/D$.  Rewriting \eqref{erc:gain-rate} gives
\[
 a_u=D_u-c_u^\partial-
       \sum_{v\sim u,\ v\text{ core}} h(D_v).
\]
Every still-uninfected frozen neighbour makes equal contributions in
the two copies.  The infected predecessor and $B_u$ infected outward
neighbours make at most $(1+B_u)/2$ difference, because
$h(D)\in[1/2,1)$ for a core.  For each union promotion after infection,
a paired child contributes at most $1/2$ difference; a single child
contributes $h(D)<1$.  Consequently, at every time,
\begin{equation}\label{erc:local-discrepancy}
 |a_u^+-a_u^-|
 \le Z_u+R_u+\frac{1+B_u}{2}.
\end{equation}
This bound includes arbitrary repeated feedback through the predecessor
and through all already infected or charged branches.

An unpaired leaf gain changes $Z_u$ by one in absolute value; its total
rate is $|a_u^+-a_u^-|$.  A union promotion increments $R_u$ by one
at rate $\max(b_u^+,b_u^-)\le1$.  For an uninfected frozen outward
neighbour $v$, the only unequal term in its aggregate leaf-gain rate
is the contribution from $u$.  Its infection rate is therefore
\begin{equation}\label{erc:edge-rate}
 \left|\frac1{D_u^+}-\frac1{D_u^-}\right|
 \le \frac{Z_u}{d(d+Z_u)}.
\end{equation}
Full degrees never decrease, which proves the inequality even if the
sign of the degree difference changes.  In particular the birth rate
of $B_u$ is at most $(d-1)Z_u/[d(d+Z_u)]\le1$.

It follows that
\[
 Q_u=Z_u+R_u+B_u
\]
has unit upward jumps with total predictable rate at most
\begin{equation}\label{erc:reservoir-rate}
 Z_u+R_u+\frac{1+B_u}{2}+1+1
 \le Q_u+\frac52.
\end{equation}
Downward jumps can only improve the comparison.  The absence of a
simultaneous promotion and neighbour infection is important here:
it follows from the aggregate-channel coupling, not from coupling
the original source clocks.

Conditionally at infection, $Q_u$ is thus dominated by a pure birth
process $Y$ started at $1$, with transition rate
\begin{equation}\label{erc:dominating-chain}
 n\longrightarrow n+1\quad\text{at rate }n+\frac52.
\end{equation}
For example, a common Poisson measure constructs this domination,
since the upward rate bound is increasing in the current state;
downward jumps of $Q_u$ preserve the order.  The comparison is uniform
in the whole prior history and all states of the attached branches.

\subsection{An exact rational reproduction certificate}

For $d\ge2$ and $s>0$ define
\begin{equation}\label{erc:mass}
 m_d(s)=\frac{d-1}{d}\,
 \mathbb E\int_0^\infty e^{-st}\frac{Y(t)}{d+Y(t)}\,dt,
 \qquad m(s)=\sup_{d\ge2}m_d(s).
\end{equation}
The compensator of the infections from $u$, together with
\eqref{erc:edge-rate} and the increasing function $z/(d+z)$, gives
\begin{equation}\label{erc:offspring}
 \mathbb E\left[
   \sum_{v:\,u\to v}e^{-s(\tau_v-\tau_u)}
      \,\middle|\,\mathcal F_{\tau_u}\right]
 \le m_d(s)\le m(s).
\end{equation}
No independence of different infected cores is asserted or needed.

\begin{lemma}[A rational discount below the original exponent]
\label{erc:certificate}
For every $s\ge s_0:=11/20$,
\begin{equation}\label{erc:subcritical-mass}
 m(s)\le\frac{19}{20}<1.
\end{equation}
\end{lemma}
\begin{proof}
It suffices to use $s=s_0$, by monotonicity in the discount.
Let $w_n$ be the expected discounted occupation time of state $n$ for
\eqref{erc:dominating-chain}.  The exponential holding times give
\[
 w_1=\frac{20}{81},\qquad
 w_{n+1}=w_n\frac{20n+50}{20n+81},\qquad
 \sum_{n\ge1}w_n=\frac{20}{11}.
\]
Consequently, for every integer $N\ge1$,
\begin{equation}\label{erc:rational-upper}
 m_d(s_0)
 =\frac{d-1}{d}\left(\frac{20}{11}
                  -d\sum_{n\ge1}\frac{w_n}{d+n}\right)
 \le\frac{d-1}{d}\left(\frac{20}{11}
                  -d\sum_{n=1}^N\frac{w_n}{d+n}\right).
\end{equation}
With $N=16$, exact rational comparison for $2\le d\le63$ gives a
maximum of the displayed upper bounds at $d=6$, equal to
\[
 \frac{1669075573443141901662377825}
      {1772128093509182310390306279}
 <\frac{19}{20}.
\]
The complete finite calculation, using only integer arithmetic and
rational fractions, is in \texttt{scripts/verify\_infection\_kernel.py}.

For the remaining degrees, the mean of $Y$ is
$\mathbb EY(t)=(7/2)e^t-5/2\le(7/2)e^t$.  Hence, for $d\ge64$,
\begin{align*}
 m_d(s_0)
 &\le\int_0^\infty e^{-t/2}
            \min\{1,(7/(2d))e^t\}\,dt\\
 &=4\sqrt{\frac7{2d}}-\frac7d
 \le\sqrt{\frac78}<\frac{15}{16}<\frac{19}{20}.
\end{align*}
The first inequality uses $s_0\ge1/2$ and
$Y/(d+Y)\le\min\{Y/d,1\}$; the expectation of this minimum is
bounded by the minimum of the expectation and one.  The last radical
comparison is equivalently the rational inequality
$7/8<(15/16)^2$.  This covers every integer degree.
\end{proof}

Generation by generation in the infection tree, conditional expectation
and \eqref{erc:offspring} now give
\begin{equation}\label{erc:total-infections}
 \mathbb E\sum_{u\text{ infected}}e^{-s\tau_u}
 \le\sum_{k\ge0}m(s)^k\le20
 \qquad(s\ge s_0).
\end{equation}
For a fully formal construction, first stop at finite population and
infection generation, use the compensator identities there, and then
pass to the limit by Tonelli and monotone convergence.  All finite
horizons are nonexplosive by real-population Yule domination.

\subsection{The old-state boundary gradient}

\begin{theorem}[One boundary leaf has a uniform resolvent effect]
\label{erc:gradient}
For every finite retained inert state $y$, every finite $T$, and
$s\ge11/20$, addition of one leaf at its root satisfies
\begin{equation}\label{erc:gradient-bound}
 \boxed{
 |V_{s,T}(y^+)-V_{s,T}(y)|
 \le\frac{2W_T(s)}{s(1-m(s))}
 \le\frac{40}{s}W_T(s).
 }
\end{equation}
The bound is uniform in the number of old cores, their degrees, and
their attached finite fringes.
\end{theorem}
\begin{proof}
Every eligible common core is present in both copies for all time
after its paired birth.  Those cores contribute equally to the two
core counts, whether or not they are eventually infected.  Every other
core belongs to a unique charged family: take its first ancestor born
from an eligible infected core after that core's infection.  These
families are disjoint within each marginal.  This partitions all
possible contributions to the absolute core-count difference.

At each union promotion from an infected core, at most two newborn
families are charged.  Conditionally on their birth time and the joint
present, each marginal has the correct future generator, so each
remaining discounted reward is bounded by $W_T(s)$.  This assertion
does not require independence from the other family or from the
infection exploration.  The birth time is a stopping time and every
such family starts in an admissible newborn environment.

The union-promotion rate at $u$ is at most one.  Its expected discounted
charge is therefore at most
\[
 2W_T(s)\,\mathbb E\int_{\tau_u}^T e^{-st}\,dt
 \le\frac{2W_T(s)}s\,\mathbb E e^{-s\tau_u}.
\]
Sum over infected eligible cores and use
\eqref{erc:total-infections}, or its sharper bound $1/(1-m(s))$.
Pathwise domination of the absolute count difference by the charged
family counts and Tonelli's theorem give \eqref{erc:gradient-bound}.
Every occurrence of $W_T$ has finite horizon, so the argument is not
assuming the uniform resolvent conclusion it will prove.
\end{proof}

\section{Closing the uniform-growth gap}
\label{erc:closure-section}

For an actual newborn environment, let $\tau_M$ be the first time its
external parent degree exceeds $M$, taking $\tau_M=0$ if it is already
larger or if the parent is an inert stub.  Let the hybrid agree with the actual process through
$\tau_M$, and thereafter delete parent-to-root arrivals.  Write
$A_M(t)$ for the supremum of its retained core mean.  Theorem~\ref{pup:hybrid} proves that if
$\mu<\lambda$ and $\delta=1/(M+1)<\lambda-\mu$, then
\begin{equation}\label{erc:hybrid-assumption}
 \gamma_M:=\limsup_{t\to\infty}t^{-1}\log A_M(t)<\lambda.
\end{equation}
This theorem includes arbitrary old states at the
crossing.  Only this strict gap for each sufficiently large fixed $M$
is needed below; no bound on its size uniform in $M$ is required.

\begin{proposition}[A return coefficient independent of the exit mass]
\label{erc:absorb}
For every $s\ge11/20$ and finite $T$,
\begin{equation}\label{erc:absorb-finite}
 W_T(s)\le\int_0^T e^{-st}A_M(t)\,dt
                 +\frac{40}{s^2(M+1)}W_T(s).
\end{equation}
In particular, if $s>\gamma_M$ and $40/[s^2(M+1)]<1$, then
\begin{equation}\label{erc:absorb-bound}
 W(s)\le
 \frac{\int_0^\infty e^{-st}A_M(t)\,dt}
      {1-40/[s^2(M+1)]}<\infty,
 \qquad\lambda\le s.
\end{equation}
\end{proposition}
\begin{proof}
After the crossing, the actual retained component has conditional
generator $\mathcal L_\infty+q_t\Delta_{\rm root}$, where
$0\le q_t\le\delta$.  For $f_r(y)=V_{s,r}(y)$, the finite-horizon
resolvent identity is
\[
 \partial_rf_r=C+(\mathcal L_\infty-s)f_r,\qquad f_0=0.
\]
Apply the time-dependent discounted Dynkin identity to
$e^{-st}f_{T-t}(Y_t)$ from $\tau_M\wedge T$ through $T$, with no
post-crossing term on $\{\tau_M\ge T\}$.  The reward before
$\tau_M$ and the inert continuation value at $\tau_M$ give exactly
the hybrid reward through $T$.  The remaining term is
\[
 \mathbb E_x\int_{\tau_M}^T e^{-st}q_t
             \Delta_{\rm root}V_{s,T-t}(Y_t)\,dt.
\]
The theorem above, $W_{T-t}(s)\le W_T(s)$, and
$\int_0^T e^{-st}dt\le1/s$ bound this term from above by
$40\delta W_T(s)/s^2$.  Localization is justified by finite-horizon
population moments; the resulting boundary increment is already
uniform in the old state.  Taking the supremum over initial newborn
environments proves \eqref{erc:absorb-finite} without interchanging a
supremum and an integral.

The coefficient is strictly less than one, so division is legitimate
at finite $T$.  Monotone convergence then proves finiteness of $W$.
Finally each $F_x(t)$ is nondecreasing, and hence
$F_x(t)\le se^{st}\int_0^\infty e^{-su}F_x(u)du\le se^{st}W(s)$.
Taking the supremum over $x$ gives $\lambda\le s$.
\end{proof}

\begin{theorem}[Equality of the original and uniform exponents]
\label{erc:equality}
The original friend-tree mean exponent and the uniform newborn mean
exponent agree:
\begin{equation}\label{erc:lambda-mu}
 \boxed{\lambda=\mu.}
\end{equation}
\end{theorem}
\begin{proof}
The existing comparison gives $\mu\le\lambda$, and the certified lower
bound on the original exponent gives $\mu\ge0.565>11/20$.
Suppose $\mu<\lambda$.  Choose a fixed integer $M\ge132$ large enough
that $1/(M+1)<\lambda-\mu$.  Then \eqref{erc:hybrid-assumption} holds.
Choose
\[
 \max\{11/20,\gamma_M\}<s<\lambda.
\]
The elementary rational inequality
$40/[(11/20)^2\cdot133]=16000/16093<1$ ensures
$40/[s^2(M+1)]<1$.  Proposition~\ref{erc:absorb} therefore yields
$\lambda\le s$, contradicting $s<\lambda$.
\end{proof}

